% -------------------------------------------------------------------------
% WBS ID: RES-PHY-SCL
% Deliverable: Dimensional Analysis, Scaling and Flow Regimes
% Status: In progress
% Evidence status: Local high-Reynolds-number tsunami-impact regime established
% Review status: Not reviewed
% Last updated: 2026-07-19
% Dependencies: RES-MOD-EQS, RES-MOD-SRC
% Downstream interfaces: RES-NUM-SPEC, RES-VER-MET
% -------------------------------------------------------------------------

\section{RES-PHY-SCL --- Dimensional Analysis, Scaling and Flow Regimes}
\label{sec:res-phy-scl}

\subsection{Purpose and Scope}
\label{sec:res-phy-scl-purpose}

% TODO: Define what this Deliverable covers and explicitly excludes.

This Deliverable records the dimensional and regime arguments that
support the local three-dimensional URANS--VOF tsunami--barrier model.
It links the high-Reynolds-number turbulence assumption to the physical
outputs required by the project.

\subsection{Research Questions}
\label{sec:res-phy-scl-research-questions}

% TODO: State the questions that must be answered before the Deliverable
% can be accepted.

\subsection{Terminology and Definitions}
\label{sec:res-phy-scl-definitions}

% TODO: Define the technical terminology, symbols, quantities and
% classifications used in this Deliverable.

\subsection{Literature Evidence}
\label{sec:res-phy-scl-literature-evidence}

% TODO: Record evidence from peer-reviewed papers, authoritative datasets,
% standards, technical reports and primary documentation.
%
% Do not create a detached literature-summary list. Explain how each source
% supports, contradicts or limits a project decision.

\textcite{qin2018comparison} is retained as evidence that local tsunami
impact can require three-dimensional velocity and nonhydrostatic
pressure resolution. \textcite{kozelkov2022turbulence} is retained as
evidence that turbulence effects become material during breaking,
run-up, overtopping and barrier loading. Source roles: `3D-RANSVOF`
and `3D-SST`. Project conclusion: the local model shall be treated as a
high-Reynolds-number turbulent flow model rather than a laminar
diagnostic calculation. Limitation: dimensionless regime arguments do
not replace validation of pressure, force and overtopping outputs.

\subsection{Governing Relations and Quantitative Definitions}
\label{sec:res-phy-scl-governing-relations}

% TODO: Record relevant equations, dimensionless groups, empirical models,
% extraction rules or quantitative definitions.
%
% Use align environments for displayed equations.
% Every displayed equation must have a unique \label{eq:...}.
% Do not place punctuation after displayed equations.

\subsection{Physical Mechanisms}
\label{sec:res-phy-scl-physical-mechanisms}

% TODO: Complete this RES-PHY Domain-specific subsection.

\subsection{Characteristic Spatial and Temporal Scales}
\label{sec:res-phy-scl-characteristic-spatial-temporal-scales}

Characteristic scales shall be recorded separately for the regional
event, the corridor extract and the local barrier domain:

\begin{itemize}
  \item regional wavelength, source footprint and propagation time from
    the Tōhoku source to the target coast;
  \item corridor width $W$, offshore/pre-impact length
    $L_{\mathrm{pre}}$, inland length $L_{\mathrm{inland}}$ and
    sponge/relaxation widths;
  \item local flow depth $h$, representative velocity $U$, barrier
    height, barrier width, gap spacing, crest length and wake length.
\end{itemize}

The corridor dimensions are physical and numerical scale choices. They
shall be derived from source footprint, severe-impact observations,
defence geometry, bathymetry/topography coverage and relaxation-zone
requirements, not from visual convenience.

\subsection{Relevant Observables}
\label{sec:res-phy-scl-relevant-observables}

% TODO: Complete this RES-PHY Domain-specific subsection.

\subsection{Controlling Dimensionless Parameters}
\label{sec:res-phy-scl-controlling-dimensionless-parameters}

% TODO: Complete this RES-PHY Domain-specific subsection.

The local flow regime shall be characterised using Reynolds, Froude and
Weber numbers:

\begin{align}
  \mathrm{Re}
  &=
  \frac{
    U L
  }{
    \nu_{\mathrm{w}}
  }
  \label{eq:res-phy-scl-reynolds-number}
  \\
  \mathrm{Fr}
  &=
  \frac{
    U
  }{
    \sqrt{
      g h
    }
  }
  \label{eq:res-phy-scl-froude-number}
  \\
  \mathrm{We}
  &=
  \frac{
    \rho_{\mathrm{w}} U^2 L
  }{
    \sigma
  }
  \label{eq:res-phy-scl-weber-number}
\end{align}

Here $U$ is a representative incident or local velocity scale, $L$ is a
local barrier or flow-depth scale, $h$ is local water depth,
$\nu_{\mathrm{w}}$ is water kinematic viscosity and $\sigma$ is
surface tension.

\subsection{Physical Regimes and Transition Conditions}
\label{sec:res-phy-scl-physical-regimes-transition-conditions}

% TODO: Complete this RES-PHY Domain-specific subsection.

The baseline local regime is turbulent, gravity-dominated,
free-surface flow with large inertial loading. High Reynolds number
supports use of a turbulence closure. Froude number identifies
supercritical or near-critical run-up, overtopping and bore-like
conditions. Large Weber number supports treating surface tension as a
minor tsunami-scale force except in small entrained droplets or
aeration details, which are outside the baseline.

The Kamaishi and Sendai corridors represent different physical regimes
within the same event. Kamaishi is retained for ria/bay and defence-
interaction scaling; Sendai is retained for coastal-plain inundation
scaling. Source role: supports active two-corridor comparison.
Supporting sources: \textcite{mori2012nationwide} for contrasting
Sanriku and Sendai inundation/run-up behaviour, and
\textcite{usgs2011tohokuEvent} for the event source point used in
trajectory construction. Limitation: final corridor bearings and
dimensions remain GIS/data products owned by `RES-DAT-SCN`.

\subsection{Implications for Model Validity}
\label{sec:res-phy-scl-model-validity-implications}

% TODO: Complete this RES-PHY Domain-specific subsection.

The high-Reynolds-number regime supports URANS $k$--$\omega$ SST as the
operational local model and LES as a higher-fidelity comparator. The
same regime also prevents free-surface height alone from validating the
model: velocity, pressure, force, impulse, overtopping and wake
quantities must be checked separately.

The corridor scale also constrains the validity of one-way replay.
If reflection, bay resonance, barrier blocking or coastal confinement
propagates materially back through the local inlet, the simultaneous
coupling option in `RES-MOD-MUL` shall be tested.

\subsection{Candidate Approaches}
\label{sec:res-phy-scl-candidate-approaches}

% TODO: Define the credible candidate approaches considered.

\subsection{Critical Comparison}
\label{sec:res-phy-scl-critical-comparison}

% TODO: Compare candidates using physically and project-relevant criteria.
% Do not select an approach solely on popularity or implementation ease.

\subsection{Assumptions and Applicability}
\label{sec:res-phy-scl-assumptions}

% TODO: State assumptions and the conditions under which they are valid.

\subsection{Limitations and Failure Conditions}
\label{sec:res-phy-scl-limitations}

% TODO: State where the evidence, model or selected approach ceases to be
% reliable or applicable.

\subsection{Project-Specific Conclusions}
\label{sec:res-phy-scl-conclusions}

% TODO: Convert the evidence into conclusions specific to the Studentship
% project. Do not merely restate literature findings.

\subsection{Selected Approach and Rationale}
\label{sec:res-phy-scl-selected-approach}

% TODO: Record the selected approach only after sufficient evidence exists.
% State the alternatives rejected and the reasons for rejection.

The selected regime interpretation is:

\begin{quote}
  local tsunami--barrier interaction is treated as high-Reynolds-number
  turbulent, gravity-dominated, three-dimensional free-surface flow
  requiring URANS turbulence closure for operational screening
\end{quote}

The selected corridor interpretation is:

\begin{quote}
  use constant-width baseline corridors for Kamaishi and Sendai until
  GIS evidence shows that a physical bay, coastline, river or
  bathymetric constraint requires a non-rectangular footprint
\end{quote}

\subsection{Unresolved Decisions}
\label{sec:res-phy-scl-unresolved-decisions}

% TODO: Record unresolved questions, required evidence and decision owners.

\subsection{Requirements and Handoffs}
\label{sec:res-phy-scl-requirements}

% TODO: State requirements passed to other Research Domains, Software,
% verification activities or later execution work.

`RES-MOD-SRC` shall retain the turbulence closure. `RES-NUM-SPEC`
shall solve the turbulence variables after pressure correction.
`RES-VER-MET` shall require validation metrics beyond free-surface
agreement.

`RES-DAT-SCN` shall provide the corridor dimensions and final bearings.
`RES-VER-TST` shall test whether the corridor and interface placement
contaminate run-up, overtopping, force, moment or energy metrics.

\subsection{Deliverable Completion Record}
\label{sec:res-phy-scl-completion-record}

% TODO: Record whether the Deliverable is:
%
% - Not started
% - In progress
% - Draft complete
% - Reviewed
% - Accepted
%
% Acceptance requires:
%
% - the research questions to be answered;
% - claims to be supported by appropriate evidence;
% - assumptions and limitations to be explicit;
% - the project-specific conclusion to be stated;
% - downstream requirements to be traceable;
% - unresolved decisions to be closed or formally deferred.
